% Segunda Lei de Kepler.
% 11/06/2003 at 21:00h, Penha, Faro
% Written by Tordar

% O foco da elipse fica em sqrt(a^2 - b^2)

\documentclass[12pt,a4paper]{article}

\pagestyle{empty}

\usepackage{graphicx}
\usepackage{pst-3dplot}

\begin{document}

\begin{pspicture}(-3,-2.5)(3,4.25)
\pstThreeDCoor[linewidth=2pt,linecolor=white,xMin=0,xMax=8,yMin=0,yMax=5,zMin=0,zMax=5,drawing=false]
\psset{linewidth=0.5mm,linecolor=black,linestyle=dashed}
\pstThreeDEllipse(0,0,0)(5,0,0)(0,4,0)
\psset{linewidth=1mm,linecolor=red,linestyle=solid}
\pstThreeDLine[linewidth=1.5mm,arrows=->](3,0,0)(0,4,0)
\pstThreeDLine[linewidth=1.5mm,arrows=->](3,0,0)(-5,0,0)
\pstThreeDLine[linewidth=1.5mm,arrows=->](0,4,0)(-5,0,0)
\psset{linewidth=1mm,linecolor=blue}
\pstThreeDLine[linewidth=2mm,arrows=->](3,0,0)(3,0,5)
\psset{dotsize=4mm,linecolor=black}
\pstThreeDDot(3,0,0)
\psset{dotsize=2mm,linecolor=black}
\pstThreeDDot(0,4,0)
\pstPlanePut[plane=xz](3.2,0,5.5){\scalebox{1.5}{$\Delta\vS=\half{1}\mathbf{r}\times\Delta\mathbf{r}$}}
\pstPlanePut[plane=xz](2.5,2,0){\scalebox{1.5}{$\mathbf{r}$}}
\pstPlanePut[plane=xz](-5.3,0,0){\scalebox{1.5}{$\mathbf{r}+\Delta\mathbf{r}$}}
\pstPlanePut[plane=xz](-2.7,2,0){\scalebox{1.5}{$\Delta\mathbf{r}$}}
\pstPlanePut[plane=xz](4,0,0){\scalebox{1.2}{$M$}}
\pstPlanePut[plane=xz](0,4.5,0){\scalebox{1.2}{$m$}}
\end{pspicture}

\end{document}
